\chapter{Literature Review}
\label{ch:literature}

\section{Overview}
\label{sec:lit-overview}

Nexis is a systems-integration project: its contribution lies in composing
established, well-engineered tools into a closed loop that none of them forms
on its own. Intellectual honesty therefore requires a clear account of what
each underlying tool already does, and where Nexis merely depends on it versus
where Nexis extends it. This section surveys the relevant landscape in that
spirit.

\textbf{LLM coding assistants.} GitHub Copilot~\cite{copilot} is the most
widely deployed AI pair-programming tool; it completes code and, in its agent
modes, drafts multi-file changes inside an editor session. It operates
entirely within the developer's authoring context: it never observes a
production incident, never validates its own output against a failing system,
and never deploys anything. Nexis's Backend agent occupies the same
``generate a minimal code change'' niche, but is invoked by a fault-detection
pipeline rather than a keystroke, and its output is machine-validated before
any human sees it.

\textbf{Durable workflow orchestration.} Temporal~\cite{temporal} provides
durable, replayable workflow execution with automatic retry and state
persistence. Nexis uses Temporal as-is, as a dependency: the entire nine-agent
RecoveryPipeline is a Temporal workflow, which is why every agent step
survives process restarts and is recorded as a durable ActivityEvent. The
workflow's engineering semantics -- role routing, contract-violation
detection, severity classification -- are Nexis's own logic layered on top.

\textbf{Graph databases for dependency analysis.} Neo4j~\cite{neo4j} is a
native graph database. Nexis stores its code-dependency graph in Neo4j and
the Pathfinder agent walks it from a fault symptom to candidate root-cause
nodes, forwarding hop distance and node degree as evidence. Neo4j is a
storage and traversal dependency; the evidence-ranking performed on the
traversal results is Nexis code.

\textbf{Causal inference tooling.} DoWhy~\cite{dowhy} is Microsoft
Research's library for explicit causal modelling with refutation tests. It is
the methodological reference point for Nexis's causal-inference sidecar, but
it must be stated plainly that Nexis does \emph{not} fit a DoWhy structural
causal model: no interventional production data exists to fit one against.
The sidecar instead implements a documented proxy -- ranking candidates by
evidence-chain specificity, graph-structural proximity, textual overlap, and
a scenario prior, with a transparent per-component confidence breakdown.

\textbf{Observability standards.} OpenTelemetry~\cite{otel} standardises
traces, metrics, and logs. Nexis adopts it as its instrumentation layer
(exported to Prometheus, Loki, and Tempo behind Grafana). It is a
dependency; Nexis adds nothing to the standard itself.

\textbf{GitOps deployment.} Argo CD~\cite{argocd} reconciles a cluster
against a Git repository and exposes Sync and Rollback operations. Nexis's
DevOps agent emits Argo CD manifests, and the post-deploy rollback probe
calls Argo CD's real Rollback client when a policy-gated SLO breach is
detected -- integration, not extension.

\textbf{Container runtimes.} Docker~\cite{docker} is the target API for both
the Validator's sandbox (network-isolated, read-only, capability-dropped
shadow execution of every candidate patch) and the preview-deploy engine.
Podman~\cite{podman} was substituted in the evaluation environment as a
Docker-API-compatible runtime and behaved identically; this substitution is
reported honestly rather than hidden.

\textbf{Property-based testing.} Hypothesis~\cite{hypothesis} generates test
inputs from declared properties. Nexis's QA agent auto-generates
pytest/Hypothesis suites from the incident and the Backend diff -- a novel
\emph{use} of the library, with Hypothesis itself unmodified. A background QA
loop replays stored suites and raises a fresh incident on a pass-to-fail
regression.

\textbf{Data lineage.} OpenLineage~\cite{openlineage} defines an open
specification for run-level lineage events. Nexis's Data Engineer agent emits
OpenLineage-compatible RunEvents for every migration it proposes or applies;
Nexis conforms to the specification rather than extending it.

\textbf{The incident-response ecosystem.} Sentry~\cite{sentry},
Datadog~\cite{datadog}, and PagerDuty~\cite{pagerduty} respectively capture
application errors, aggregate observability data, and route alerts to on-call
humans. All three are integrated into Nexis as webhook-driven incident
\emph{sources} feeding the same Sentinel pipeline. The Slack
API~\cite{slackapi} is used in the opposite direction: the Approval Gate
pushes a decision into Slack, where an engineer can approve or reject a
validated patch with one click. Supporting platform dependencies -- WorkOS
for SSO~\cite{workos}, Stripe for billing~\cite{stripe}, PostgreSQL
row-level security for tenant isolation~\cite{pgrls}, pgvector for
knowledge-base retrieval~\cite{pgvector}, and Ollama~\cite{ollama} as the
local LLM provider used for this report's live evidence -- are used as
documented, without modification. Table~\ref{tab:lit-tools} summarises the
dependency-versus-contribution boundary.

\begin{table}[htbp]
\caption{Position of surveyed tools relative to Nexis.}
\label{tab:lit-tools}
\centering
\small
\begin{tabular}{@{}lll@{}}
\toprule
\textbf{Tool} & \textbf{What it provides} & \textbf{Relation to Nexis} \\
\midrule
Copilot~\cite{copilot} & Editor-time code generation & Comparable to one agent (Backend) only \\
Temporal~\cite{temporal} & Durable workflow execution & Dependency; hosts the RecoveryPipeline \\
Neo4j~\cite{neo4j} & Graph storage and traversal & Dependency; Pathfinder walks it \\
DoWhy~\cite{dowhy} & Structural causal models & Reference point; proxy method used instead \\
OpenTelemetry~\cite{otel} & Telemetry standard & Dependency; instrumentation layer \\
Argo CD~\cite{argocd} & GitOps sync/rollback & Integrated by DevOps agent and rollback probe \\
Docker~\cite{docker}/Podman~\cite{podman} & Container runtime & Dependency; sandbox and deploy engine \\
Hypothesis~\cite{hypothesis} & Property-based testing & Novel use: QA agent generates suites \\
OpenLineage~\cite{openlineage} & Lineage specification & Conformance by Data Engineer agent \\
Sentry/Datadog/PagerDuty~\cite{sentry,datadog,pagerduty} & Error/alert capture & Webhook incident sources \\
Slack API~\cite{slackapi} & Team messaging & One-click approval channel \\
\bottomrule
\end{tabular}
\end{table}

\section{Research Gap}
\label{sec:lit-gap}

Each tool above is strong within its boundary, and several are routinely
combined in industry -- Sentry alerting into PagerDuty, an on-call engineer
opening an editor with Copilot, Argo CD deploying the eventual merge. What no
single tool, and no combination in common use, provides is a system that
(a)~models a software engineering \emph{team} with explicit role
specialisation, and (b)~closes the loop from automatic fault detection to a
deployed, verified fix, with human judgement concentrated at one
severity-routed decision point.

The gap is visible at every seam of the conventional tool chain. Copilot
generates code but never sees a production incident; it has no notion of a
fault, a blast radius, or a deployment~\cite{copilot}. Observability
platforms such as Datadog and Sentry detect and display faults with great
fidelity but never generate a candidate fix, let alone validate
one~\cite{datadog,sentry}. PagerDuty routes a human to the problem; the
repair itself remains entirely manual~\cite{pagerduty}. Argo CD deploys
whatever a repository contains but authors nothing~\cite{argocd}. Temporal
can durably orchestrate any of these steps but carries no engineering
semantics of its own~\cite{temporal}. In every case the loop is broken at
the same place: between ``a fault is known'' and ``a validated fix exists,''
a human performs unaided diagnosis, authoring, testing, review, and release.

Nexis closes precisely that span, end to end. Sentinel detects a fault using
a statistical-process-control rule (an EWMA baseline with a
mean-plus-three-sigma breach condition), not a fixed threshold. Pathfinder
walks the real Neo4j dependency graph to candidate root causes, which the
causal-inference sidecar ranks with a transparent confidence breakdown. The
Synthesiser classifies the scenario and selects which Execution Team agents
are actually needed -- and the workflow enforces that routing, visibly
skipping unselected agents. The Architect produces a structured plan whose
\code{affected\_files} contract is mechanically checked against the Backend
agent's actual diff; a violation forces HIGH severity so the change can never
auto-deploy unreviewed. QA generates property-based tests from the incident
and the diff~\cite{hypothesis}; the Validator executes patch and tests in a
locked-down container sandbox, so no unvalidated patch ever reaches a human.
Only then does the single human decision point engage -- and only when
severity warrants it: HIGH-severity changes require explicit review,
MEDIUM-severity changes auto-approve after a 120-second window unless
overridden, and trivial LOW-severity changes pass immediately. The decision
itself is richer than a binary gate: the engineer may Approve, Reject, or
Modify the diff before it ships, with every terminal decision persisted for
future preference learning. On approval, a real GitHub pull request is opened
and a post-deploy rollback probe watches for regression~\cite{argocd}.

Two further properties distinguish this from a scripted pipeline of existing
tools. First, role specialisation is structural, not cosmetic: five Execution
Team agents (Architect, Backend, QA, DevOps, Data Engineer) and four
Self-Healing Loop agents (Sentinel, Pathfinder, Synthesiser, Validator) have
distinct inputs, outputs, and contracts, and the Architect's plan constrains
the others. Second, the loop is closed over the platform's own
infrastructure: when Nexis's preview-deploy engine fails to build or
health-check a repository, that failure is inserted as a first-class incident
into the very same pipeline, and ``write a working Dockerfile'' becomes an
ordinary repair task. To the author's knowledge, no tool in common use today
combines these properties in a single system.

\section{Conclusion}
\label{sec:lit-conclusion}

The surveyed landscape supplies every ingredient Nexis needs -- durable
orchestration~\cite{temporal}, graph traversal~\cite{neo4j}, sandboxed
execution~\cite{docker,podman}, property-based testing~\cite{hypothesis},
GitOps delivery~\cite{argocd}, and a mature incident-capture
ecosystem~\cite{sentry,datadog,pagerduty} -- but leaves the loop between
detection and repair open, to be closed by a human. Nexis's contribution is
not any single new algorithm; it is the demonstrated composition of these
tools into a role-specialised multi-agent team that detects, diagnoses,
repairs, validates, and deploys autonomously, engaging human judgement once,
at a severity-routed gate. The remaining chapters describe how that
composition was specified, designed, implemented, and tested.
